\begin{lemma}
Under the two-bite law, condition on the red base graph $G_R$ and the red row map $\rho$. For any fixed set $T$ of original vertices and any fixed target set $S$ of blue coordinates, the conditional probability that all vertices in $T$ are assigned blue coordinates in $S$ is bounded by $(|S|/M)^{|T|}$, where $M$ is the base size.
\end{lemma}
\begin{proof}
Write a two-bite sample as \(\omega=(G'_R,G_B,\phi)\), where
\(\phi:\operatorname{Fin} n\hookrightarrow \operatorname{Fin} m\times
\operatorname{Fin} m\).  Let
\[
C=\{\omega:G'_R=G_R\text{ and }(\phi(v))_1=\rho(v)\text{ for every }v\}
\]
be the conditioning cell, and let
\[
E=\{\omega:(\phi(v))_2\in S\text{ for every }v\in T\}
\]
be the event whose conditional probability is being bounded.  By
\(\noderef{TwoBiteConditionalProbability}\), if
\(\Pr(C)=0\), then \(\Pr(E\mid C)\) is defined to be \(0\).  The right hand
side \((|S|/M)^{|T|}\) is nonnegative: if \(M>0\), then \(|S|/M\ge 0\), and
if \(M=0\), then \(S\subseteq\operatorname{Fin}0\), so \(|S|=0\) and the
real quotient is \(0\).  Thus the desired inequality holds in the zero-mass
case.

Assume from now on that \(\Pr(C)\ne 0\).  Let
\[
X_{\rm all}
=\{\phi:\operatorname{Fin} n\hookrightarrow
\operatorname{Fin} m\times\operatorname{Fin} m:
(\phi(v))_1=\rho(v)\text{ for all }v\},
\]
and let
\[
X_{\rm good}
=\{\phi\in X_{\rm all}:(\phi(v))_2\in S\text{ for every }v\in T\}.
\]
Set \(a=|X_{\rm all}|\) and \(g=|X_{\rm good}|\).  Also set
\[
w_R=\operatorname{GnpGraphWeight}_p(G_R),\qquad
u=\operatorname{UniformInjectionWeight}(n,m),\qquad
B=\sum_{H}\operatorname{GnpGraphWeight}_p(H),
\]
where \(H\) ranges over all blue base graphs on \(\operatorname{Fin} m\).
The definitions \(\noderef{TwoBiteEventProbability}\) and
\(\noderef{TwoBiteSampleWeight}\) give the conditioning mass as
\[
\Pr(C)
=\sum_{H}\sum_{\phi\in X_{\rm all}}w_R\,
  \operatorname{GnpGraphWeight}_p(H)\,u
=w_R\,B\,u\,a.
\]
Likewise,
\[
\Pr(E\cap C)
=\sum_{H}\sum_{\phi\in X_{\rm good}}w_R\,
  \operatorname{GnpGraphWeight}_p(H)\,u
=w_R\,B\,u\,g.
\]
These equalities are just finite rearrangements of the sample sum: the red
graph is fixed to be \(G_R\), the blue graph \(H\) is unconstrained, and the
injection variable is exactly the set \(X_{\rm all}\), or \(X_{\rm good}\)
after imposing \(E\).  The factors \(w_R\), \(u\), and
\(\operatorname{GnpGraphWeight}_p(H)\) do not depend on the injection
\(\phi\), so they factor as shown.

The assumption \(\Pr(C)\ne0\) rules out every zero case that would make the
conditioning ill-defined.  In particular \(w_R\,B\,u\ne0\) and
\(a\ne0\).  Hence the fixed red graph has nonzero weight in the relevant
cell, the blue-graph total factor is nonzero, the global injection weight is
nonzero, and there is at least one admissible injection with first coordinate
\(\rho\).  Therefore the common factor \(w_R\,B\,u\) cancels in the
conditional probability:
\[
\Pr(E\mid C)
=\frac{\Pr(E\cap C)}{\Pr(C)}
=\frac{w_R\,B\,u\,g}{w_R\,B\,u\,a}
=\frac{g}{a}.
\]
If the red row map \(\rho\) has an impossible fiber, for example more than
\(M\) vertices assigned to one red row, then \(X_{\rm all}\) is empty and
\(a=0\); this is one of the zero-mass cases already disposed of above.

It remains only to bound the finite fraction \(g/a\).  This is exactly the
content of \(\noderef{OppositeProjectionRowInjectionCounting}\): for the
nonempty admissible-injection set \(X_{\rm all}\), the fraction of injections
whose vertices in \(T\) have blue coordinate in \(S\) is at most
\((|S|/M)^{|T|}\).  Combining this counting inequality with the identity
\(\Pr(E\mid C)=g/a\) proves the stated conditional bound.
\end{proof}
